% ===================================================================
%  तक्ता क्रमांक १५ — बाजार. --- खाद्यपदार्थ.
%
%  Traduction, faite en 2026, en marathi d'aujourd'hui, sur l'ido de
%  texto/io. Regles de la colonne : voir 10-takta-01.tex. Une seule
%  suite d'alineas numerotes : les cles ne portent pas d'indice de
%  scene.
%
%  LE TABLEAU LE PLUS FACILE DE LA COLONNE, ET C'EST LA SYNTAXE QUI
%  LE DIT. Il enumere, il n'accroche pas : la postposition ne mord
%  que sur « X de Y », et un etal n'est qu'une liste. Le tableau 13,
%  trois fois plus court, avait coute neuf inversions ; celui-ci en
%  coute presque aucune. C'est exactement ce que la colonne
%  cantonaise avait constate au meme tableau, dans une langue sans
%  aucun rapport.
%
%  ET LE MARCHE EST MARATHE PRESQUE MOT POUR MOT, parce qu'un marche
%  ne contient que des choses qui existent partout :
%
%    बटाटा (61)   la pomme de terre   कांदा (—)    l'oignon
%    लसूण (101)   l'ail               गाजर         la carotte
%    मुळा         le radis            वाटाणे       les petits pois
%    घेवडा (63)   les haricots        मसूर (59)    les lentilles
%    तांदूळ (58)  le riz              अंजीर (64)   les figues
%    बदाम (56)    les amandes         केळे (46)    la banane
%    दुधी भोपळा (19) la courge        खरबूज (18)   le melon
%    फुलकोबी (15) le chou-fleur       कोबी (17)    le chou
%    चुका         l'oseille           आळीव        le cresson
%
%  ET LA MER DONNE SES MOTS DE CONKAN : कोळीण la poissonniere — du
%  nom des Koli, les pecheurs de la cote —, वाम l'anguille, शिंपले
%  les moules, कालव l'huitre, खेकडा le crabe, कोळंबी la crevette,
%  शेवंड la langouste. Le tableau 10 avait donne les mots du port ;
%  celui-ci donne ceux de l'etal.
%
%  ET « खाटीक » NOMME LE BOUCHER PAR SA CASTE, comme les neuf metiers
%  du tableau 14. La marchande de tripes, le garcon boucher et
%  l'ecailliere en derivent tous.
%
%  CE QUI NE POUSSE PAS ICI GARDE SON NOM D'AILLEURS : l'artichaut,
%  le camembert, le gruyere, le hareng, le rhum, le kirsch, le
%  faisan. Le jambon aussi, faute d'un mot marathe pour une cuisse
%  de porc salee. Meme regle qu'au marron chaud du tableau 7.
% ===================================================================

%%K t15-apar-1 apar
\VUsaut{4.0mm}
\VUtitre{15.0pt}{60}{तक्ता क्रमांक १५}
\VUsaut{3.0mm}

%%K t15-tit-1 sub
\VUpk{11.4pt}{40}{बाजार. --- खाद्यपदार्थ.}
\VUsaut{3.0mm}

%%K t15-01-1 p
१. --- आपल्या खाण्यासाठी लागणारा माल आपल्याला पुरवणारे बहुतेक व्यापारी
\VUgras{छपऱ्या बाजाराच्या} \VUgras{दालनाखाली} \textsuperscript{(1)} किंवा शेजारच्या रस्त्यांवर
जमतात.

%%K t15-02-1 p
२. --- \VUgras{डुकराचा खाटीक} \textsuperscript{(2)}, जो \VUgras{हॅम} \textsuperscript{(3)},
\VUgras{रक्ताची सॉसेज} \textsuperscript{(4)}, \VUgras{सॉसेज} \textsuperscript{(5)},
\VUgras{आतड्याची सॉसेज} \textsuperscript{(6)} आणि \VUgras{खारवलेली चरबी} \textsuperscript{(7)}
विकतो, तो \VUgras{लोणी आणि चीज विकणारीच्या} \textsuperscript{(8)} शेजारी उभा असतो, जिची मांडणी
लाल कवचाच्या \VUgras{डच चीजने} \textsuperscript{(9)}, \VUgras{कामँबेर चीजने}
\textsuperscript{(10)}, \VUgras{ग्रुयेरच्या चकत्यांनी} \textsuperscript{(11)} आणि ताज्या
\VUgras{लोण्याच्या गोळ्यांनी} \textsuperscript{(12)} भरलेली असते.

%%K t15-03-1 p
३. --- तिच्यासमोर एक \VUgras{भाजीवाली} \textsuperscript{(13)} आपल्या गिऱ्हाइकिणींना सुंदर
\VUgras{टोमॅटो} \textsuperscript{(14)}, \VUgras{फुलकोबी} \textsuperscript{(15)},
\VUgras{कोशिंबिरीची पाने} \textsuperscript{(16)}, \VUgras{कोबी} \textsuperscript{(17)},
\VUgras{दुधी भोपळे} \textsuperscript{(19)}, \VUgras{खरबुजे} \textsuperscript{(18)} आणि
\VUgras{अळंबीसुद्धा} \textsuperscript{(20)} देऊ करते. पण एका \VUgras{बिअरवाल्याने} आपली
\VUgras{मालाची गाडी} \textsuperscript{(21)} — \VUgras{बिअरच्या पिंपांनी}
\textsuperscript{(22)} भरलेली — नेमकी तिच्या मांडणीसमोर उभी केल्यामुळे तिच्या गिऱ्हाइकांना जवळ
येता येत नाही. एक माळीण तिच्याकडे \VUgras{आर्टिचोकची} \textsuperscript{(23)} टोपली घेऊन येते.

%%K t15-tit-2 sub
\VUsaut{4.0mm}
\VUpk{10.2pt}{40}{विक्री आणि खरेदी.}
\VUsaut{3.0mm}

%%K t15-04-1 p
४. --- बाजारात मोठी लगबग आहे, कारण \VUgras{लिलावाची} \textsuperscript{(24)} वेळ झाली.
खरेदीसाठी तिथे येणाऱ्या \VUgras{गृहिणी} \textsuperscript{(26)} जाताजाता
\VUgras{आतडी विकणारीच्या} \textsuperscript{(25)} दुकानासमोर थांबतात, \VUgras{पोटातली आतडी}
किंवा \VUgras{वासराचे डोके} घ्यायला.

%%K t15-05-1 p
५. --- एक लठ्ठ \VUgras{आचारी} \textsuperscript{(27)} \VUgras{कोळिणीकडे} \textsuperscript{(28)}
फार सुंदर \VUgras{टुन्याचा} भाव करतो आहे, जी आपल्या मर्जीप्रमाणे आपले भाव वाढवते. तिच्याकडे
\VUgras{वाम}, \VUgras{सोल मासे}, \VUgras{ट्राउट} आणि \VUgras{कार्प} मोठ्या प्रमाणात आले आहेत.
तिचा \VUgras{कॉड} आणि तिचे \VUgras{शिंपले} फार ताजे आहेत.

%%K t15-tit-3 sub
\VUsaut{4.0mm}
\VUpk{10.2pt}{40}{स्वस्त आणि महाग माल.}
\VUsaut{3.0mm}

%%K t15-06-1 p
६. --- तिच्याशेजारी बसलेली \VUgras{कोंबडीवाली} \textsuperscript{(29)} फार प्रामाणिक आहे.
\VUgras{टर्की}, \VUgras{हंस}, \VUgras{बदक} किंवा साधे \VUgras{कोंबडीचे पिल्लू} विकताना ती
रास्त भावापेक्षा जास्त मागत नाही.

%%K t15-07-1 p
७. --- चौकाच्या कोपऱ्यावर, पहिल्या मजल्यावर, थोड्या दिवसांपूर्वी एक \VUgras{कापडविक्रेता}
\textsuperscript{(30)} बसला आहे, ज्याच्याकडे खरेदी करणाऱ्या बायकांना \VUgras{मेजपोसांची},
\VUgras{रुमालांची}, \VUgras{अपऱ्यांची} आणि \VUgras{फडक्यांची} फार सुंदर निवड मिळेल याची नेहमी
खात्री असते. त्याच मजल्यावर एका \VUgras{सुऱ्या करणाऱ्याने} \textsuperscript{(31)} आपली
कार्यशाळा थाटली आहे. तो दालनातल्या विक्रेतीणींच्या \VUgras{सुऱ्या} धारदार करतो आणि माळ्यांसाठी
सर्वात कठीण \VUgras{पोलादाचे} \VUgras{कोयते} बनवतो.

%%K t15-tit-4 sub
\VUsaut{4.0mm}
\VUpk{10.2pt}{40}{किराणा.}
\VUsaut{3.0mm}

%%K t15-08-1 p
८. --- त्याच्या कार्यशाळेखाली थोड्या दिवसांपूर्वी अन्नपदार्थांचे एक मोठे दुकान उघडले आहे. ते
आता पूर्वीचे साधे आणि किरकोळ \VUgras{किराणा दुकान} \textsuperscript{(32)} राहिलेले नाही, जे
फक्त रोजच्या वापरातला माल विकत असे — \VUgras{चहा} \textsuperscript{(33)}, \VUgras{चॉकलेट}
\textsuperscript{(34)}, \VUgras{साखरेचे ठोकळे} \textsuperscript{(37)} आणि \VUgras{साबण}
\textsuperscript{(38)}. तिथे कोणताही माल मिळतो : \VUgras{खारी बिस्किटे},
\VUgras{डबाबंद पदार्थ} \textsuperscript{(35)}, \VUgras{मद्ये} \textsuperscript{(36)},
\VUgras{रम}, \VUgras{कोन्याक}, \VUgras{किर्श}, \VUgras{बडीशेपेचे मद्य} आणि \VUgras{क्युरासाओ}.
तिथे शिकारही मिळते : \VUgras{ससा} \textsuperscript{(39)}, \VUgras{थ्रश}
\textsuperscript{(40)}, \VUgras{तित्तिर} \textsuperscript{(41)}, \VUgras{लावे}
\textsuperscript{(42)}, \VUgras{रानबदक} \textsuperscript{(43)} किंवा \VUgras{फेझंट}
\textsuperscript{(44)} — तितक्याच सहजपणे जितक्या सहजपणे \VUgras{केळ्यांसारखी}
\textsuperscript{(46)} आणि \VUgras{अननसासारखी} परदेशी फळे.

%%K t15-09-1 p
९. --- किराणावाल्याचा एक फार सौजन्यशील \VUgras{नोकर} \textsuperscript{(45)} हवे ते द्यायला
तयार असतो : बेदाण्यांचा किंवा त्या फळांचा \VUgras{मुरंबा} \textsuperscript{(47)},
\VUgras{चरबी} \textsuperscript{(48)}, \VUgras{लोणच्याच्या काकड्या} \textsuperscript{(49)}
किंवा खारवलेले \VUgras{सार्डिन} \textsuperscript{(50)}.

%%K t15-09-2 p
हे \VUgras{गिऱ्हाइकिणींसाठी} \textsuperscript{(51)} फार सोयीचे आहे, ज्यांना रोजच्या
मालाशेजारीच सर्वात दुर्मीळ नवी पिके आणि सर्वात चवदार फळे सापडतात : रसाळ \VUgras{नासपती}
\textsuperscript{(52)}, \VUgras{आलुबुखार} \textsuperscript{(53)}, \VUgras{द्राक्षे}
\textsuperscript{(54)}, \VUgras{पीच} \textsuperscript{(55)}, \VUgras{बदाम}
\textsuperscript{(56)} आणि \VUgras{अक्रोड} \textsuperscript{(57)}.

%%K t15-tit-5 sub
\VUsaut{4.0mm}
\VUpk{10.2pt}{40}{गिऱ्हाईक.}
\VUsaut{3.0mm}

%%K t15-10-1 p
१०. --- \VUgras{तांदूळ} \textsuperscript{(58)} आणि \VUgras{मसूर} \textsuperscript{(59)} धुरी
दिलेल्या \VUgras{हेरिंगच्या} \textsuperscript{(60)} पेटीशेजारी आहेत; \VUgras{बटाटा}
\textsuperscript{(61)} आणि \VUgras{घेवडा} \textsuperscript{(63)} \VUgras{सफरचंदांच्या}
\textsuperscript{(62)}, \VUgras{अंजिरांच्या} \textsuperscript{(64)},
\VUgras{सुक्या आलुबुखारांच्या} \textsuperscript{(65)} आणि \VUgras{हेझलनटांच्या}
\textsuperscript{(66)} शेजारी आहेत. त्या दुकानात ताजी \VUgras{अंडी} \textsuperscript{(67)} आणि
\VUgras{ऑलिव्हही} \textsuperscript{(68)} मिळतात; म्हणूनच \VUgras{टोपल्या}
\textsuperscript{(70)} आणि \VUgras{जाळीच्या पिशव्या} \textsuperscript{(73)} भराभर भरून जातात.

%%K t15-11-1 p
११. --- त्या उत्तम पेढीच्या \VUgras{कॉफीने} \textsuperscript{(72)} \VUgras{नोकराणींमध्ये}
\textsuperscript{(69)} आणि \VUgras{स्वयंपाकिणींमध्ये} रास्त नाव कमावले आहे, कारण म्हातारा
\VUgras{किराणावाला} \textsuperscript{(71)} स्वतःच ती अत्यंत काळजीपूर्वक भाजतो.

%%K t15-tit-6 sub
\VUsaut{4.0mm}
\VUpk{10.2pt}{40}{देखावे.}
\VUsaut{3.0mm}

%%K t15-12-1 p
१२. --- सगळेच काही खरेदीसाठी बाजारात येत नाहीत. दालनाकडे जाणाऱ्या गल्लीतल्या चित्रमय वर्दळीत
अत्यंत निरनिराळी माणसे दिसतात. एका मनमिळाऊ \VUgras{खाटकाच्या नोकराशेजारी}
\textsuperscript{(75)}, जो आपल्यापुढे \VUgras{ढकलगाडी} \textsuperscript{(74)} ढकलतो, हे बघा
रजेवरचे दोन शिपाई, आपला झळाळता गणवेश दाखवायला फार अभिमानाने उभे : निळा \VUgras{डोलमान} आणि
\VUgras{तुऱ्याची टोपी} घातलेला \VUgras{हुसार} \textsuperscript{(76)}, आणि फुगीर बाह्यांचा व
डोक्यावर \VUgras{फेझ} घातलेला \VUgras{झुआव नाईक}.

%%K t15-13-1 p
१३. --- \VUgras{पाववाला} \textsuperscript{(80)}, जो \VUgras{पावाच्या दुकानाच्या}
\textsuperscript{(78)} उंबरठ्यावर उभा आहे, त्याचे आपल्या \VUgras{पावाचे}
\textsuperscript{(79)} पीठ नुकतेच मळून झाले आहे, आणि दर्शनी काचेत असलेले
\VUgras{चंद्रकोरी क्रोसाँ} \textsuperscript{(81)}, \VUgras{गोल पाव} \textsuperscript{(82)} आणि
\VUgras{चकतीचे पाव} \textsuperscript{(83)} त्याने भट्टीतून बाहेर काढले आहेत.

%%K t15-14-1 p
१४. --- पावाच्या दुकानावर एक कुशल \VUgras{शिवणकाम करणारी} \textsuperscript{(84)} राहते, जी
कित्येक \VUgras{भरतकाम करणाऱ्या} \textsuperscript{(85)} बायका कामाला ठेवते. तिच्याशेजारी एक
\VUgras{इस्त्रीवाली} \textsuperscript{(86)} राहते, जी \VUgras{कपडे} \textsuperscript{(88)}
धुऊन वाळवल्यावर त्यांना खळ लावते आणि \VUgras{इस्त्रीने} \textsuperscript{(87)} इस्त्री करते.

%%K t15-tit-7 sub
\VUsaut{4.0mm}
\VUpk{10.2pt}{40}{मांस, मासे आणि भाज्या.}
\VUsaut{3.0mm}

%%K t15-15-1 p
१५. --- तळमजल्यावरच्या \VUgras{खाटीकखान्यात} \textsuperscript{(89)} सगळ्या प्रकारचे
\VUgras{मांस} विकतात — \VUgras{मेंढीचे मांस} \textsuperscript{(90)}, \VUgras{गोमांस}
\textsuperscript{(94)} किंवा \VUgras{वासराचे मांस} \textsuperscript{(92)} —
\VUgras{मेंढीच्या मांड्या}, \VUgras{बीफस्टेक}, \VUgras{भाजायचे तुकडे} आणि \VUgras{चॉप} म्हणून.
\VUgras{खाटकाचा नोकर} \textsuperscript{(93)} ते सगळे मांस कापून वेगळे करतो आणि
\VUgras{मज्जेची हाडे} \textsuperscript{(96)} बाजूला ठेवतो. खाटीकखान्याच्या दारात एक
\VUgras{कालवे विकणारी} \textsuperscript{(97)} बसते, जी \VUgras{कालवे} \textsuperscript{(98)}
विकते. हवे असल्यास ती ती उघडूनही देते.

%%K t15-16-1 p
१६. --- हे सगळे व्यापारी परवाना किंवा जागेचा कर भरतात; म्हणूनच \VUgras{पोलिस}
\textsuperscript{(100)} त्या बिचाऱ्या \VUgras{भाजी विकणाऱ्या फेरीवाल्यांचा}
\textsuperscript{(99)} निर्दयपणे पाठलाग करतो, ज्या आपल्या हातगाड्यांवरून \VUgras{गाजरे},
\VUgras{मुळे}, \VUgras{सलगम}, \VUgras{लीक} किंवा \VUgras{ॲस्पॅरॅगसचे जुडगे},
\VUgras{ताजे वाटाणे}, \VUgras{पालक}, \VUgras{चुका}, \VUgras{आळीव} आणि \VUgras{पार्सली} विकत
फिरून त्यांच्याशी स्पर्धा करतात.

%%K t15-tit-8 sub
\VUsaut{4.0mm}
\VUpk{10.2pt}{40}{पोलिसापासून सावध!}
\VUsaut{3.0mm}

%%K t15-17-1 p
१७. --- \VUgras{लसूण} \textsuperscript{(101)} आणि \VUgras{कांदे} विकणाऱ्या पोराला चांगले माहीत
आहे की त्यालाही आपला माल बाजाराजवळ विकता येणार नाही. तो पळ काढतो, आणि धावताना \VUgras{शेवंड} व
\VUgras{लॉबस्टर} विकणाऱ्या \VUgras{विक्रेत्याची} \textsuperscript{(102)} \VUgras{खेकड्यांची},
\VUgras{गोड्या पाण्यातल्या खेकड्यांची} आणि \VUgras{कोळंबीची} टोपली पाडायचा धोका पत्करतो.

%%K t15-18-1 p
१८. --- सगळे धावपळ करतात आणि गोंगाट करतात; पण त्यामुळे फिरत्या \VUgras{काचकाम करणाऱ्याचा}
\textsuperscript{(103)} आवाज किंवा \VUgras{कोळसेवाल्याची} \textsuperscript{(104)} हाक स्पष्ट
ऐकू यायला अडथळा येत नाही — तो \VUgras{कोळशाची पोती} \textsuperscript{(105)} पोहोचवायला काही
क्षणांसाठी आपली गाडी सोडून गेला आहे.

\VUsaut{6.0mm}
\VUfilet{22mm}
